\chapter{Problem Definition and Motivation}
\label{ch:problem_statement}

\section{Problem Statement}

The automated classification of skin lesions from dermatoscopic images remains a challenging computer vision and machine learning problem characterized by several interconnected challenges:

\subsection{The Multi-Class Classification Challenge}

The task involves discriminating among seven distinct skin lesion categories (Melanoma, Nevus, Basal Cell Carcinoma, Actinic Keratosis, Benign Keratosis, Dermatofibroma, and Vascular Lesions) based on dermatoscopic images. This is fundamentally a multi-class classification problem with the following specific challenges:

\begin{enumerate}
    \item \textbf{High Inter-class Similarity}: Many lesion types exhibit remarkable visual similarity in color, texture, and morphological features. For example, melanomas can appear superficially similar to large benign nevi, requiring careful analysis of subtle discriminative features.
    
    \item \textbf{Intra-class Variation}: Individual lesion types exhibit high visual variability depending on location, stage of development, patient demographics, and imaging conditions. A single melanoma can present with diverse morphological characteristics.
    
    \item \textbf{Severe Class Imbalance}: The HAM10000 dataset exhibits significant class imbalance, with melanoma comprising only 2.6\% of samples (326 images) compared to nevus at 31.6\% (3165 images). This imbalance poses algorithmic challenges and can lead to biased models that underperform on critical minority classes.
    
    \item \textbf{Limited Annotated Data}: High-quality annotated dermatoscopic datasets remain relatively scarce compared to other computer vision domains. The HAM10000 dataset contains only 10,015 images, limiting the direct application of large-scale deep learning approaches.
\end{enumerate}

\section{Clinical Significance}

The clinical importance of accurate skin lesion classification cannot be overstated:

\subsection{Diagnostic Accuracy}

Melanoma diagnosis presents particular challenges due to its serious prognosis and mortality impact. Delayed diagnosis significantly reduces treatment efficacy and patient survival rates. Conversely, false positive diagnoses lead to unnecessary biopsies, patient anxiety, and healthcare costs.

Studies demonstrate that even experienced dermatologists exhibit significant diagnostic variability, with accuracy ranging from 75-95\% depending on expertise and lesion characteristics \cite{esteva2019dermatologist}. Computational systems that consistently achieve high accuracy could augment clinical decision-making and improve diagnostic confidence.

\subsection{Healthcare Access and Equity}

The global shortage of qualified dermatologists creates a critical bottleneck preventing timely diagnosis for millions of patients in developing regions. Automated screening systems could dramatically improve access to early detection and appropriate referral to expert dermatologists.

\subsection{Economic Impact}

Automated diagnostic systems could reduce healthcare costs by:
\begin{itemize}
    \item Minimizing unnecessary biopsies and follow-up appointments
    \item Enabling efficient resource allocation by prioritizing high-risk cases
    \item Reducing diagnostic delays that would otherwise increase treatment complexity and costs
\end{itemize}

\section{Technical Challenges}

Beyond the inherent classification difficulty, several technical challenges complicate the development of robust solutions:

\subsection{Data Imbalance Mitigation}

Traditional machine learning often assumes balanced classes. The severe imbalance in skin lesion datasets requires sophisticated techniques including:

\begin{itemize}
    \item Oversampling of minority classes
    \item Weighted loss functions accounting for class frequencies
    \item Threshold adjustment for optimal operating points
    \item Stratified sampling strategies
\end{itemize}

\subsection{Transfer Learning and Domain Adaptation}

While pre-training on ImageNet provides strong initialization, the domain gap between natural images and dermatoscopic images can limit direct transferability. The model must learn domain-specific features including:

\begin{itemize}
    \item Dermoscopic patterns (patterns of pigment distribution)
    \item Subtle color variations imperceptible to non-experts
    \item Texture characteristics specific to different lesion types
\end{itemize}

\subsection{Limited Batch Size Sensitivity}

Medical datasets typically support only modest batch sizes. Understanding how architectural design impacts performance under batch size constraints is crucial for practical deployment.

\subsection{Computational Efficiency}

Clinical deployment requires computationally efficient models that can run on resource-constrained hardware while maintaining accuracy. The trade-off between model complexity and computational requirements must be carefully managed.

\section{Architectural Considerations}

Recent architectural innovations offer potential solutions to these challenges:

\subsection{Vision Transformer Architecture}

Vision Transformers address limitations of convolution-based approaches:

\begin{enumerate}
    \item \textbf{Global Receptive Fields}: Attention mechanisms capture long-range dependencies more effectively than convolutional receptive fields, enabling the model to simultaneously consider morphological features across different image regions.
    
    \item \textbf{Flexibility in Feature Extraction}: Transformers learn specialized attention patterns for different lesion types, potentially capturing subtle discriminative features more effectively than fixed convolutional kernels.
    
    \item \textbf{Hierarchical Representations}: Recent transformer variants including Swin Transformers provide hierarchical multi-scale representations aligned with modern CNN architectures.
\end{enumerate}

\subsection{Swin Transformer Specific Advantages}

The Swin Transformer architecture offers particular advantages for medical imaging:

\begin{itemize}
    \item \textbf{Computational Efficiency}: Shifted window attention mechanisms reduce computational complexity to linear compared to quadratic for standard attention, enabling practical deployment.
    
    \item \textbf{Hierarchical Design}: Multi-stage architecture with progressively reduced spatial resolution aligns with standard CNN conventions, facilitating transfer learning and feature pyramid construction.
    
    \item \textbf{Inductive Biases}: Explicit incorporation of locality and translational equivariance through shifting windows provides beneficial inductive biases for vision tasks.
    
    \item \textbf{Recent Success in Medical Imaging}: Recent applications to medical image segmentation and classification demonstrate Swin Transformer effectiveness on domain-specific tasks.
\end{itemize}

\section{Research Questions}

This thesis addresses the following research questions:

\begin{enumerate}
    \item How does the Swin Transformer architecture compare to traditional CNN approaches on skin lesion classification, and what are the specific sources of performance gains?
    
    \item What architectural modifications and training strategies are necessary to optimize Swin Transformer performance for medical imaging tasks with limited data?
    
    \item How effectively can shift window attention mechanisms capture discriminative features specific to different skin lesion types?
    
    \item What is the practical feasibility of deploying Swin Transformer models in clinical settings with computational and latency constraints?
    
    \item How does performance vary across different lesion types, and are there specific classes for which transformer architectures provide particular advantages?
\end{enumerate}

\section{Thesis Scope}

This thesis focuses specifically on:

\begin{itemize}
    \item Multi-class skin lesion classification on the HAM10000 dataset
    \item Swin Transformer architecture application and optimization
    \item Comprehensive performance evaluation with clinical-relevant metrics
    \item Ablation studies examining architectural component contributions
    \item Practical deployment considerations and recommendations
\end{itemize}

Out of scope are:

\begin{itemize}
    \item Clinical validation in real dermatology practices
    \item Longitudinal follow-up studies
    \item Other skin imaging modalities (3D imaging, polarized light microscopy)
    \item Explanation methods for clinical interpretability (though briefly discussed)
\end{itemize}

This problem definition establishes the foundation for the comprehensive study presented in subsequent chapters.
